\documentclass[11pt]{article}
\usepackage[margin=1in]{geometry}
\usepackage{amsmath,amssymb,amsthm}
\usepackage{graphicx}
\usepackage{hyperref}

\newtheorem{theorem}{Theorem}
\newtheorem{remark}{Remark}

\newcommand{\PD}{\mathbb{P}}
\newcommand{\R}{\mathbb{R}}
\newcommand{\tr}{\operatorname{tr}}
\newcommand{\diag}{\operatorname{diag}}

\title{A Holder-Datum Counterexample to Strict Hilbert Contractivity}
\author{Run \texttt{fa\_banach\_001}, lane 15}
\date{2026-07-03}

\begin{document}
\maketitle

\section*{Source and Open Question}

The source paper is Max Fathi, Samuel Sra, and Melanie Weber,
\emph{Computing Brascamp-Lieb Constants through the lens of Thompson Geometry},
arXiv:2208.05013.  In Section 6, after introducing the normalized Picard map
\[
  \widetilde G(X)=\frac{G(X)}{\|G(X)\|},
\]
for an arbitrary matrix norm, the authors ask whether $\widetilde G$ is
strictly contractive with respect to Hilbert's projective metric.  The relevant
source passage is reproduced below.

\begin{center}
\includegraphics[width=0.9\linewidth]{figures/open_problem_crop.png}
\end{center}

\section*{Proof Intuition}

Hilbert's projective metric ignores positive scalar multiples.  Therefore the
normalization in $\widetilde G=G/\|G\|$ cannot create projective contraction if
the underlying map $G$ is already the identity on a feasible datum.  The
Holder datum with repeated identity maps gives exactly this situation.

\section*{Definitions}

For positive definite matrices $X,Y\in \PD_d$, Hilbert's projective metric is
\[
  \delta_H(X,Y)=\log\bigl(M(X/Y)M(Y/X)\bigr),
  \qquad
  M(X/Y)=\inf\{\lambda>0:X\preceq \lambda Y\}.
\]
The source Picard map is
\[
  G(X)=\left(\sum_{j=1}^m w_j A_j(A_j^\ast X A_j)^{-1}A_j^\ast\right)^{-1}.
\]
Strict contractivity means
\[
  \delta_H(\widetilde G(X),\widetilde G(Y))<\delta_H(X,Y)
\]
for all projectively distinct positive definite matrices $X,Y$ in the domain.

\section*{Counterexample}

\begin{theorem}
The unqualified assertion that $\widetilde G$ is strictly contractive in
Hilbert's projective metric is false for feasible Brascamp-Lieb data.
\end{theorem}

\begin{proof}
Fix any dimension $d\ge 2$ and consider the Holder Brascamp-Lieb datum
\[
  m=2,\qquad A_1=A_2=I_d,\qquad w_1=w_2=\frac12.
\]
This is feasible: it is the standard Holder datum from the source paper, with
$\sum_j w_j=1$.  Equivalently, Tao's feasibility conditions give
\[
  d=\sum_{j=1}^2 w_j d
\]
and, for every subspace $H\subseteq \R^d$,
\[
  \dim H=\sum_{j=1}^2 w_j\dim(I_dH).
\]

For this datum, the Picard map is exactly the identity:
\[
\begin{aligned}
  G(X)
  &=\left(\frac12 I_d(I_d^\ast X I_d)^{-1}I_d^\ast
          +\frac12 I_d(I_d^\ast X I_d)^{-1}I_d^\ast\right)^{-1}  \\
  &=\left(\frac12 X^{-1}+\frac12 X^{-1}\right)^{-1}
   =X.
\end{aligned}
\]
Hence
\[
  \widetilde G(X)=\frac{X}{\|X\|}
\]
for every positive definite $X$ and every matrix norm used in the
normalization.

Hilbert's projective metric is invariant under positive scalar rescaling:
for all $\alpha,\beta>0$,
\[
  \delta_H(\alpha X,\beta Y)=\delta_H(X,Y).
\]
Indeed $M(\alpha X/\beta Y)=(\alpha/\beta)M(X/Y)$ and
$M(\beta Y/\alpha X)=(\beta/\alpha)M(Y/X)$, so the two scalar factors cancel.
Therefore
\[
  \delta_H(\widetilde G(X),\widetilde G(Y))
  =\delta_H(X,Y)
\]
for every pair of projectively distinct positive definite matrices $X,Y$.

For an explicit pair in dimension $2$, take
\[
  X=\diag(2,1),\qquad Y=I_2.
\]
Then $M(X/Y)=2$ and $M(Y/X)=1$, so $\delta_H(X,Y)=\log 2>0$, and equality of
distances proves that $\widetilde G$ is not strictly contractive.
\end{proof}

\begin{remark}
This counterexample is intentionally narrow.  The Holder datum above is
feasible but not simple: every nonzero proper subspace $H$ is critical, since
\[
  \dim H=\sum_{j=1}^2 w_j\dim(I_dH).
\]
Thus the example refutes the literal general feasible-datum formulation of the
Section 6 question, but it does not settle a possible intended version
restricted to simple Brascamp-Lieb data.
\end{remark}

\section*{Verification Notes}

The script \texttt{code/verify\_holder\_identity.py} checks the displayed
$2\times 2$ witness using the diagonal formula
\[
  \delta_H(\diag(x_1,x_2),\diag(y_1,y_2))
  = \log\frac{\max_i x_i/y_i}{\min_i x_i/y_i}.
\]
It prints equality of the pre- and post-iteration distances, both equal to
$\log 2$.

\section*{Novelty and Limitations}

Local run indexes were searched for arXiv:2208.05013, the paper title,
Brascamp-Lieb, Thompson geometry, Hilbert projective metric, and strict
contractivity.  A bounded web search on 2026-07-03 for the exact
open-question phrases, arXiv id, and close variants found no separate answer
source.  Because the example is elementary, novelty confidence is only
moderate.

Human review should check whether the intended source question silently
assumes simple data.  Under that stronger reading, this packet should be read
as a literal-scope caveat, not as a solution of the simple-data strict
contractivity problem.

\begin{thebibliography}{9}
\bibitem{FathiSraWeber}
Max Fathi, Samuel Sra, and Melanie Weber.
\emph{Computing Brascamp-Lieb Constants through the lens of Thompson Geometry}.
arXiv:2208.05013.
\end{thebibliography}

\end{document}
